\section*{الفصل \ref{ch:b3:renal-osmoregulation} --- فيزيولوجيا الكلية وتنظيم الأسمولية}

\begin{solution}{exo:b3:renal-osmoregulation:1}
\omterm{def:b3:renal-osmoregulation:nephron}{الكبيبة}: ترشّح البلازما ناقصةً البروتينات. \omterm{def:b3:renal-osmoregulation:nephron}{والنبيب القريب}:
يعيد امتصاص ثلثي الملح والماء، وكل الغلوكوز والأحماض
الأمينية، والبيكربونات. والفرع النازل من \omterm{def:b3:renal-osmoregulation:nephron}{عروة هنلي}: يفقد
الماء إلى اللب. والفرع الصاعد: يضخ NaCl خارجًا، وهو غير نفوذ
للماء. والنبيب البعيد: يضبط الصوديوم (\omterm{def:b3:renal-osmoregulation:controls}{بالألدوستيرون})، ويفرز
البوتاسيوم. \omterm{def:b3:renal-osmoregulation:nephron}{والقناة الجامعة}: إعادة امتصاص الماء تحت
الفازوبريسين، ومعالجة البروتونات والبولة.
\end{solution}

\begin{solution}{exo:b3:renal-osmoregulation:2}
التصفية: حجم البلازما المصفّى من المادة في وحدة الزمن، أي
$UV/P$. وهنا $100\times 1.5/2 = \qty{75}{mL/min}$، دون معدل الرشح
البالغ \qty{120}{mL/min}: فالمادة تُرشَّح ويُعاد امتصاص جزء
منها (نحو \qty{38}{\%} من المقدار المرشَّح).
\end{solution}

\begin{solution}{exo:b3:renal-osmoregulation:3}
سمكة الماء العذب أملح من وسطها: فيدخل الماء عبر الخياشيم
بالتناضح، فالشرب لا يزيد إلا في الطوفان؛ وهي تطرح بولًا
غزيرًا مخففًا وتلتقط خلاياها الكلوريدية في الخياشيم
Na$^{+}$ وCl$^{-}$ التقاطًا فاعلًا من الماء. وسمكة البحر أقل
ملوحة من البحر: فتفقد الماء عبر الخياشيم، وتشرب ماء البحر
لتعوّضه، وتضخ الخلايا الكلوريدية الملح المبتلع إلى الخارج.
\end{solution}

\begin{solution}{exo:b3:renal-osmoregulation:4}
\omterm{prop:b3:renal-osmoregulation:countercurrent}{الأثر المفرد} هو فرق \qty{200}{mOsm/L} الذي يبقيه الفرع الصاعد
الغليظ بين الخلال وسائله هو بضخ NaCl خارجًا من غير أن يدع
ماءً يتبعه. والفرع الصاعد يجب أن يكون غير نفوذ للماء، وإلا
تبع الماءُ الملحَ وزال الفرق؛ والفرع النازل يجب أن يكون
نفوذًا للماء حتى يكون السائل البالغ المنعطف قد غُلّظ أصلًا
بالخلال وتعمل كل مرحلة على سائل أملح من سابقتها. ولا تقدر
شعبتان متساويتا النفوذية على المضاعفة.
\end{solution}

\begin{solution}{exo:b3:renal-osmoregulation:5}
الطرف الوارد: $60 - 18 - 28 = \qty{14}{mmHg}$؛ والطرف الصادر: $60 - 18 -
35 = \qty{7}{mmHg}$. ومع ترشيح الماء تتركز البروتينات الباقية
ويرتفع الضغط الجرمي على امتداد الشعيرة؛ فإذا بلغ
\qty{42}{mmHg} صار الضغط الصافي صفرًا وتوقف الرشح قبل النهاية
(توازن الرشح). ويتوقف معدل الرشح عندئذ على الجريان البلازمي:
فالجريان الأسرع يركّز البروتينات أبطأ، فيستمر الرشح على
امتداد أطول ويرتفع معدل الرشح مع الجريان.
\end{solution}

\begin{solution}{exo:b3:renal-osmoregulation:6}
الحمل المرشَّح: $\qty{40}{mL/min}\times\qty{0.14}{mmol/mL} =
\qty{5.6}{mmol/min}$. والطرح: $70\times 1 = \qty{0.07}{mmol/min}$.
والمعاد امتصاصه: $5.53/5.6 = \qty{98.75}{\%}$. وتصفية الصوديوم: $70\times
1/140 = \qty{0.5}{mL/min}$.
\end{solution}

\begin{solution}{exo:b3:renal-osmoregulation:7}
$900/1200 = \qty{0.75}{L}$ في اليوم. وعند \qty{400}{mOsm/L}: $900/400 =
\qty{2.25}{L}$ من البول، فيجب أن يُشرب نحو $2.25 + 0.9 = \qty{3.15}{L}$
(ناقصًا ماء الطعام): فالمريض الذي يخفق آلية تغليظه عليه أن
يشرب، وهو في خطر إن لم يقدر.
\end{solution}

\begin{solution}{exo:b3:renal-osmoregulation:8}
(أ) البيلة التفهة: \qtyrange{10}{15}{L} من البول عند
\qtyrange{50}{100}{mOsm/L}، وصوديوم البلازما مرتفع كلما تأخر
الشرب، وعطش لا ينقطع. (ب) الفازوبريسين غير الملائم: بول شحيح
مغلَّظ، واحتباس ماء، وصوديوم بلازمي منخفض بالتخفيف، وعطش غير
مرتفع مع استمرار الشرب --- وتشوش واختلاجات حين يهبط الصوديوم
كثيرًا. (ج) فرط \omterm{def:b3:renal-osmoregulation:controls}{الألدوستيرون}: احتباس الصوديوم مع الماء،
وتوسع الحجم، وارتفاع ضغط الدم، وفقد البوتاسيوم (انخفاض
K$^{+}$ في البلازما)، وصوديوم البلازما قريب من السويّ إذ يتبع
الماء الملح وتفلت الكلية جزئيًا؛ وحجم البول سويّ. (د) حمية
بلا ملح: يرتفع الرينين \omterm{def:b3:renal-osmoregulation:controls}{والألدوستيرون} في غضون يوم، ويهبط
صوديوم البول إلى بضعة ميلي مولات في اليوم، ويُحفظ صوديوم
البلازما، ويهبط الحجم والضغط قليلًا.
\end{solution}

\begin{solution}{exo:b3:renal-osmoregulation:9}
البول المخفف: التصفية \omterm{def:b3:renal-osmoregulation:problem}{الأسمولية} $100\times 6/300 = \qty{2}{mL/min}$؛
\omterm{prop:b3:renal-osmoregulation:water}{وتصفية الماء الحر} $6 - 2 = +\qty{4}{mL/min}$: فالكلية تطرح
ماءً خاليًا من المذاب، كما بعد حمل مائي والفازوبريسين مطفأ.
والبول المغلَّظ: التصفية \omterm{def:b3:renal-osmoregulation:problem}{الأسمولية} $1000\times 0.6/300 =
\qty{2}{mL/min}$؛ \omterm{prop:b3:renal-osmoregulation:water}{وتصفية الماء الحر} $0.6 - 2 = -\qty{1.4}{mL/min}$:
أي إن \qty{1.4}{mL/min} من الماء الحر تُعاد إلى الجسم ---
تجفاف، والفازوبريسين مرتفع.
\end{solution}

\begin{solution}{exo:b3:renal-osmoregulation:10}
رقّم المراحل من القشرة؛ وليكن الخلال عند المرحلة $k$ هو
$I_{k}$. فالسائل النازل عند $k$ يساوي $I_{k}$؛ والسائل الصاعد
عند $k$ هو $I_{k} - 200$؛ والسائل الصاعد عند $k$ هو ما غادر
المنعطف ومرّ بالمراحل $n, n-1, \dots, k+1$. وفي الحالة
المستقرة يغادر السائل الداخل عند \qty{300}{} قمةَ الفرع
الصاعد عند $I_{1} - 200$، ويتجاوز خلال كل مرحلة ما فوقه بما
لا يزيد على \omterm{prop:b3:renal-osmoregulation:countercurrent}{الأثر المفرد}، فيكون $I_{k} \le 300 + 200k$ مع
المساواة حين تعمل كل مرحلة بأثر كامل: فيبلغ المنعطف $300 + 200n$.
فالتدرج محدود إذن بطول العروة (بعدد المراحل) مضروبًا في
\omterm{prop:b3:renal-osmoregulation:countercurrent}{الأثر المفرد}، لا بقدرة ضخ أي خلية. و$n$ هو طول العروة مقسومًا
على المسافة التي يتوازن عليها السائل مع الخلال: فالعرى
المجاورة للّب الطويلة، والعرى الطويلة جدًا عند قوارض الصحراء،
تعطي $n$ كبيرًا؛ ويحدّ الأثرَ قدرةُ المضخة وغسل المذاب
بالأوعية المستقيمة.
\end{solution}

\begin{solution}{exo:b3:renal-osmoregulation:11}
\omterm{def:b3:renal-osmoregulation:comparative}{الماء الاستقلابي}: $5\times 0.8 = \qty{4}{g}$ من النشاء تعطي $4\times 0.6
= \qty{2.4}{g}$. والفواقد: البول $3/5500\ \unit{L} = \qty{0.55}{mL}$؛
والتبخر \qty{1.5}{g} (وهي صافية بعد استعادة الأنف). فمجموع
الفقد \qty{2.05}{g} في مقابل كسب \qty{2.4}{g}: أي فائض
\qty{0.35}{g}، يغطي الفقد الرجيعي الصغير، وماء البذور الماصّ
للرطوبة زيادة. فهو ينجو من غير أن يشرب. ولولا استعادة الأنف
لكان فقد التنفس $1.5/0.3 = \qty{5}{g}$، ولمات في أيام قليلة.
\end{solution}

\begin{solution}{exo:b3:renal-osmoregulation:12}
في الجريان المتوازي يتوازن الدم والسائل الديالي في منتصف
الغشاء ويهبط فرق التركيز الذي يسوق الانتشار إلى الصفر:
فيغادر الدم في أحسن الأحوال بنصف بولته. وفي الجريان المعاكس
يلقى أطرى سائل ديالي أنظفَ دم ويُصان التدرج على امتداد
الغشاء كله، فيقدر الدم على المغادرة بتركيز يقارب تركيز مدخل
السائل الديالي --- وهو مبدأ \omterm{def:b3:renal-osmoregulation:nephron}{عروة هنلي} نفسه. والتصفية:
$300\times 0.7 =
\qty{210}{mL/min}$ على مدى $3\times 4 = \qty{12}{h}$ في الأسبوع، أي
$210\times 720 = \qty{151}{L}$ في الأسبوع؛ وكليتان عند
\qty{125}{mL/min} تصفّيان $125\times\num{10080} = \qty{1260}{L}$.
فالآلة تؤدي نحو ثُمن عمل الكليتين، أي \qty{15}{mL/min} بالمتوسط
الزمني --- يكفي للعيش، لا للعافية.
\end{solution}

\begin{solution}{pb:b3:renal-osmoregulation:1}
\textbf{1.} الضغط الصافي $55 - 15 - 30 = \qty{10}{mmHg}$؛ ومعدل الرشح $= 12.5
\times 10 = \qty{125}{mL/min}$.
\textbf{2.} $125\times 1440 = \qty{180}{L/d}$؛ و$180/3 = 60$ مرة.
\textbf{3.} الإنولين: $62.5\times 1/0.5 = \qty{125}{mL/min}$، مساوية
لمعدل الرشح. والكرياتينين: $1.3\times 1/0.01 = \qty{130}{mL/min}$،
أكبر قليلًا لأن \omterm{def:b3:renal-osmoregulation:nephron}{النبيب القريب} يفرز مقدارًا صغيرًا من
الكرياتينين فوق المرشَّح.
\textbf{4.} تصفية PAH $13\times 1/0.02 = \qty{650}{mL/min}$، وهي
الجريان البلازمي الكلوي؛ والجريان الدموي $650/(1 - 0.45) = \qty{1180}{mL/min}$،
أي نحو \qty{1.2}{L/min}؛ ونسبة الرشح $125/650 = 0.19$.
\textbf{5.} الضغط الصافي \qty{15}{mmHg}، ومعدل الرشح
\qty{188}{mL/min} (وأقل في الواقع، إذ يرتفع الضغط الجرمي على
امتداد الشعيرة). والأنجيوتنسين II يقبض الشُّرين الصادر ويرفع
الضغط الكبيبي؛ وحجبه يخفض الضغط ومعدل الرشح قليلًا. وفي
السكري ترشّح الكبيبات تحت ضغط عالٍ فيبلى المرشِّح؛ وخفض
الضغط يبطئ الأذية سنين.
\textbf{6.} الخلايا القدمية وحُجُبها الشقّية (أو شحنة الغشاء
القاعدي): وهي الطبقة التي تحبس الألبومين. وفقد \qty{3}{g/d}
يخفض الضغط الجرمي في البلازما؛ فيغادر السائل الشعيرات في كل
مكان وتنتفخ النسج --- وهي وذمة المتلازمة الكلائية.
\textbf{7.} المرشَّح $125\times 5 = \qty{0.625}{mmol/min}$؛ والمطروح
$250\times 1 = \qty{0.25}{mmol/min}$؛ والمعاد امتصاصه $0.375$، أي
\qty{60}{\%}؛ والتصفية $250/5 = \qty{50}{mL/min}$.
\textbf{8.} المرشَّح $180\times 0.14 = \qty{25.2}{mol/d}$ من Na$^{+}$،
أي \qty{580}{g} من الصوديوم، و\qty{1.47}{kg} ملحًا. والمطروح $100\times 1.44
= \qty{144}{mmol/d}$؛ والمعاد امتصاصه $1 - 144/\num{25200} = \qty{99.4}{\%}$.
فكان سيُفقد نحو كيلوغرام ونصف كيلوغرام من الملح في اليوم.
\textbf{9.} المرشَّح $125\times 5 = \qty{0.625}{mmol/min}$، أي
\qty{112}{mg/min}، و\qty{162}{g/d}؛ ولا يُطرح شيء، إذ هذا دون
$T_{m}$. ويبلغ الحمل المرشَّح $T_{m}$ عند $2.1/0.125 =
\qty{16.8}{mmol/L}$ (\qty{300}{mg/dL}) --- أعلى من العتبة
المرصودة البالغة نحو \qty{10}{mmol/L}، لأن النفرونات تختلف
وأضعفها ينسكب أولًا (تفلطح منحنى المعايرة).
\textbf{10.} المرشَّح $125\times 20 = \qty{2.5}{mmol/min}$؛ والمطروح
$2.5 - 2.1 = \qty{0.4}{mmol/min}$، أي \qty{72}{mg/min}، و$576$ ميلي مول أو
\qty{104}{g} في اليوم. والبول الإضافي $576/300 = \qty{1.9}{L}$.
فالغلوكوز يجر الماء إلى البول (إدرار أسموزي)، فيجف المريض
ويعطش؛ و\qty{104}{g} من الغلوكوز \qty{1.8}{MJ} تُفقد في اليوم،
والخلايا، إذ تعجز عن التقاط الغلوكوز، تحرق الشحم والبروتين:
فينقص الوزن.
\textbf{11.} المعاد امتصاصه $\approx \qty{25}{mol/d}$؛ وATP اللازم $25/3 =
\qty{8.4}{mol/d}$؛ و$8.4\times 50 = \qty{420}{kJ}$، أي نحو \qty{6}{\%}
من \qty{7}{MJ}.
\textbf{12.} لتفرز الكلية كل فضلة كان على النبيب أن يحمل
ناقلًا يتعرفها --- أي آلافًا منها، ولا ناقل لجزيء لم يُلقَ
قط. وأما ترشيح كل صغير وإعادة امتصاص المئات القليلة من
الأنواع التي يريدها الجسم فلا يحتاج إلا إلى ناقلات لما
يُعرف أنه نفيس؛ وكل ما لا يُتعرَّف يغادر تلقائيًا. والثمن
طاقة السؤال~11.
\textbf{13.} التصفية \omterm{def:b3:renal-osmoregulation:problem}{الأسمولية} $600\times 1/290 = \qty{2.07}{mL/min}$؛
\omterm{prop:b3:renal-osmoregulation:water}{وتصفية الماء الحر} $1 - 2.07 = -\qty{1.07}{mL/min}$: فالكلية تحفظ
الماء.
\textbf{14.} الأدنى $600/1200 = \qty{0.5}{L}$؛ والأقصى $600/50 =
\qty{12}{L}$؛ والحاجة اليومية عند الأدنى $0.5 + 0.9 = \qty{1.4}{L}$.
\textbf{15.} \qty{1000}{mOsm} عند \qty{1200}{mOsm/L} تحتاج إلى
\qty{0.83}{L} من البول: فالكسب الصافي \qty{0.17}{L} قبل مذاب اليوم.
وبعده: مجموع مذاب اليوم \qty{1600}{mOsm}، أي \qty{1.33}{L} من
البول، مع \qty{0.9}{L} غير محسوس، فيخرج \qty{2.23}{L} مقابل
\qty{1}{L} داخلًا --- أي عجز \qty{1.23}{L}، في مقابل
\qty{1.4}{L} بلا شرب. فلتر ماء البحر لا يحسّن الميزان إلا
بمقدار \qty{0.17}{L}.
\textbf{16.} البول عند \qty{50}{mOsm/L} لليوم: $600/50 =
\qty{12}{L}$؛ فيجب أن يُشرب \qty{12.9}{L} كل يوم، أي نحو \qty{13}{L}.
\textbf{17.} $B$ لترًا من الجعة تجلب $20B$ ميلي أسمول؛ وأقصى بول
هو $(600 + 20B)/50 = 12 + 0.4B$ لترًا، وعليه أن يحمل $B - 0.9$
لترًا: أي $B - 0.9 \le 12 + 0.4B$، ومنه $B \le \qty{21.5}{L}$. أي
نحو عشرين لترًا --- غير أنه إن لم يُؤكل شيء يُذكر هبط مذاب
اليوم إلى \qty{200}{mOsm} وهبط الحد إلى نحو \qty{8}{L}: وهي
نقص صوديوم الدم عند مدمني الشرب الذين لا يأكلون.
\textbf{18.} $298/290 - 1 = \qty{2.8}{\%}$: أي أبعد كثيرًا من
\qty{1}{\%} التي تطلق الفازوبريسين والعطش؛ وكلاهما أقصى ما
يكون. وبثبات المذاب: $290\times 42 = 298\times V$، ومنه
$V = \qty{40.9}{L}$، أي عجز قدره نحو \qty{1.1}{L}.
\textbf{19.} ستة لترات من الماء تخفف البلازما بينما يزيل العرق
ملحًا؛ والجهد والألم والكرب تطلق الفازوبريسين بصرف النظر عن
\omterm{def:b3:renal-osmoregulation:problem}{الأسمولية}، فلا تقدر الكلية على تخفيف البول بما يكفي لطرح
الزائد. فيهبط صوديوم البلازما، وتنتفخ خلايا الدماغ ---
وتلَت ذلك اختلاجات ووفيات. وكان على الكلية أن تُطفئ
الفازوبريسين وأن تخرج \qty{6}{L} من البول عند $U_{\min}$؛
وينبغي للعدّاء أن يشرب بحسب عطشه وأن يعوّض الملح.
\textbf{20.} $3/5500\ \unit{L} = \qty{0.55}{mL}$ في اليوم. وكلية
بشرية كانت ستحتاج إلى $3/1200 = \qty{2.5}{mL}$ للمذاب نفسه:
فجرذ الكنغر يستعمل خُمس الماء.
\textbf{21.} الداخل: \qty{2.4}{g}. والخارج: $1.6 + 0.3 + 0.55 = \qty{2.45}{g}$.
فالميزان متوازن في حدود \qty{0.05}{g} في اليوم، يوفرها الماء
الماصّ للرطوبة في البذور (بضع في المئة من \qty{5}{g}).
\textbf{22.} الماء $500\times 0.65 = \qty{325}{L}$؛ ورُبعه
\qty{81}{L}؛ وعند \qty{5}{L/d}، ستة عشر يومًا.
\textbf{23.} $150/800 = \qty{0.19}{L}$، أي نحو \qty{190}{mL} من
المحلول الملحي --- وهو أقل من \qty{100}{mL} من ماء البحر وماء
السمك معًا، فيكسب الطائر ماءً. وكليته، عند \qty{600}{mOsm/L}
(\qty{300}{mmol/L} من NaCl)، كانت ستحتاج إلى \qty{0.5}{L} من
البول لحمل الملح نفسه، أي ماءً أكثر مما أدخل: فالكلية التي
لا تقدر على التفوق على ماء البحر لا تدع طائرًا يشربه،
والغدة، عند ضعف تركيز البحر، تقدر.
\textbf{24.} \qty{30}{g} من الماء داخلة، فنحو \qty{30}{mL} من
البول في اليوم، عند كسر صغير من \omterm{def:b3:renal-osmoregulation:problem}{أسمولية} الدم (بضع عشرات من
mOsm/L في مقابل \num{300}). ويغادر الملح في ذلك البول
وبالانتشار عبر الخياشيم؛ ومن دون التقاط فاعل بالخلايا
الكلوريدية من ماء يحوي ميلي مولًا في اللتر، كانت السمكة
ستُستنزف في أيام.
\textbf{25.} معدل الرشح \qty{125}{mL/min}، أي \qty{180}{L} في
اليوم؛ والبول الإجباري \qty{0.5}{L} في اليوم؛ ولتر ماء البحر
يعطي صافيًا \qty{0.17}{L} بعد طرح ملحه --- أي لا شيء تقريبًا.
\end{solution}
